\newpage
\section*{LỜI NÓI ĐẦU}
\thispagestyle{empty}
Trong bối cảnh chuyển đổi số đang diễn ra mạnh mẽ ở khắp các lĩnh vực, nông nghiệp thông minh (Smart Agriculture) dần trở thành xu hướng tất yếu nhằm nâng cao năng suất và chất lượng nông sản. Tại Việt Nam, ngành chăn nuôi cũng đang có bước chuyển mình rõ rệt từ mô hình hộ gia đình nhỏ lẻ sang các trang trại tập trung, ứng dụng công nghệ cao.

Tuy nhiên, bài toán khó nhất đối với các trang trại quy mô lớn là khâu kiểm soát vi khí hậu chuồng nuôi. Những biến động thất thường về nhiệt độ, độ ẩm hay sự tích tụ khí độc (NH3, H2S, CO2) nếu không được xử lý kịp thời sẽ làm vật nuôi bị stress, sinh bệnh và gây thiệt hại kinh tế nặng nề. Rõ ràng, cách giám sát thủ công truyền thống vừa tốn sức người, vừa đứt quãng và không còn đáp ứng được yêu cầu thực tế.

Xuất phát từ nhu cầu đó, em quyết định chọn đề tài: "NGHIÊN CỨU THIẾT KẾ VÀ CHẾ TẠO HỆ THỐNG IOT GIÁM SÁT VI KHÍ HẬU CHUỒNG NUÔI GIA CẦM ỨNG DỤNG LORA VÀ ĐIỆN TOÁN BIÊN". Mục tiêu của đồ án là phát triển một hệ thống giám sát tự động, kết hợp thiết bị IoT và mạng truyền thông vô tuyến tầm xa (LoRa) để thu thập dữ liệu môi trường liên tục. Điểm cốt lõi là dữ liệu được xử lý ngay tại thiết bị biên theo hướng Điện toán biên (Edge Computing), giúp người chăn nuôi theo dõi và phản ứng nhanh với sự cố theo thời gian thực, kể cả khi đường truyền Internet bị gián đoạn.

Báo cáo này sẽ trình bày chi tiết toàn bộ quá trình thực hiện đồ án, từ bước khảo sát lý thuyết, lựa chọn linh kiện, thiết kế mạch phần cứng thu thập dữ liệu, cho đến khâu lập trình và kiểm thử hệ thống. Dù đã nỗ lực hoàn thiện, nhưng do giới hạn về thời gian và kinh nghiệm thực tiễn, đồ án chắc chắn không tránh khỏi những thiếu sót. Em rất mong nhận được những nhận xét và góp ý quý báu từ quý thầy cô để đề tài được hoàn thiện hơn.

Em xin chân thành cảm ơn!

\cleardoublepage